% GENERATED FILE. Edit canonical lessons, then run npm run generate:books.
% canonical-source-hash: fnv1a64:0af23ef29e70f7d7
% canonical-lessons: SA-C304-rasih, SA-C304-artah, SA-C304-krosati, SA-C304-vilapati, SA-C304-nihsarayati

\chapter{Trapped}
\label{ch:sa-trapped}
\hlchaptermodality{304}

\begin{chapteropening}
\textbf{By the end of this chapter:} \emph{I can say who was trapped and how they called for help.}

Complete the last lesson of chapter 304: I can say who was trapped and how they called for help.
\end{chapteropening}

\section[\texorpdfstring{\sk{राशिः} (rāśiḥ)}{राशिः (rāśiḥ)}]{\sk{राशिः} (rāśiḥ) — a heap; a pile}

\label{lesson:SA-C304-rasih}



\begin{quote}
Before the new one: say the Sanskrit for spreads through, then the Sanskrit for withers.
\end{quote}

\subsection*{You'll want to know: \sk{राशिः}}
\textbf{\sk{राशिः}} — \emph{rāśiḥ} — \textquotedblleft{}a heap; a pile\textquotedblright{}.

\begin{quote}
\sk{भग्नायाः} \sk{भित्तेः} \sk{राशिः} \sk{मार्गे} \sk{अस्ति।}
\end{quote}

\begin{quote}
\emph{bhagnāyāḥ bhitteḥ rāśiḥ mārge asti.}
\end{quote}

\begin{quote}
A heap from the broken wall lies in the road.
\end{quote}

A masculine noun.

\begin{cousinweb}
\textbf{\sk{राशिः}} (\emph{rāśiḥ}) MW gives \emph{rāśi} as a heap, mass or pile and calls its derivation doubtful, so its origin is uncertain (MW).
\end{cousinweb}

\begin{grammarlens}[title={What you've built}]
One more word for telling what happened.
\end{grammarlens}

\subsection*{Your turn}
\begin{itemize}
\raggedright
  \item \emph{Say it:} \emph{rāśiḥ}
  \item \emph{Say it:} \emph{rāśiḥ}, once more
  \item \emph{Say it:} \emph{bhagnāyāḥ bhitteḥ rāśiḥ mārge asti.} — A heap from the broken wall lies in the road.
  \item \emph{Recall:} say the Sanskrit for spreads through, then the Sanskrit for withers, then say \emph{rāśiḥ} again
\end{itemize}

\begin{tcolorbox}[breakable,colback=teal!4,colframe=teal!35!black,title={Before you move on}]
What does \sk{राशिः} mean? (\textquotedblleft{}A heap; a pile\textquotedblright{}.) Say it once more.
\end{tcolorbox}
\section[\texorpdfstring{\sk{आर्तः} (ārtaḥ)}{आर्तः (ārtaḥ)}]{\sk{आर्तः} (ārtaḥ) — in distress; afflicted}

\label{lesson:SA-C304-artah}



\begin{quote}
Before the new one: say the Sanskrit for withers, then the Sanskrit for a heap.
\end{quote}

\subsection*{You'll want to know: \sk{आर्तः}}
\textbf{\sk{आर्तः}} — \emph{ārtaḥ} — \textquotedblleft{}in distress; afflicted\textquotedblright{}.

\begin{quote}
\sk{राशेः} \sk{अधः} \sk{आर्तः} \sk{नरः} \sk{अस्ति।}
\end{quote}

\begin{quote}
\emph{rāśeḥ adhaḥ ārtaḥ naraḥ asti.}
\end{quote}

\begin{quote}
There is a man in distress under the heap.
\end{quote}

\begin{cousinweb}
\textbf{\sk{आर्तः}} (\emph{ārtaḥ}) MW gives \emph{ārta} as fallen into misfortune, afflicted, pained, and discusses its derivation without settling it, so its origin is uncertain (MW).
\end{cousinweb}

\begin{grammarlens}[title={What you've built}]
One more word for telling what happened.
\end{grammarlens}

\subsection*{Your turn}
\begin{itemize}
\raggedright
  \item \emph{Say it:} \emph{ārtaḥ}
  \item \emph{Say it:} \emph{ārtaḥ}, once more
  \item \emph{Say it:} \emph{rāśeḥ adhaḥ ārtaḥ naraḥ asti.} — There is a man in distress under the heap.
  \item \emph{Recall:} say the Sanskrit for withers, then the Sanskrit for a heap, then say \emph{ārtaḥ} again
\end{itemize}

\begin{tcolorbox}[breakable,colback=teal!4,colframe=teal!35!black,title={Before you move on}]
What does \sk{आर्तः} mean? (\textquotedblleft{}In distress; afflicted\textquotedblright{}.) Say it once more.
\end{tcolorbox}
\section[\texorpdfstring{\sk{क्रोशति} (krośati)}{क्रोशति (krośati)}]{\sk{क्रोशति} (krośati) — shouts; cries out}

\label{lesson:SA-C304-krosati}



\begin{quote}
Before the new one: say the Sanskrit for a heap, then the Sanskrit for in distress.
\end{quote}

\subsection*{You'll want to know: \sk{क्रोशति}}
\textbf{\sk{क्रोशति}} — \emph{krośati} — \textquotedblleft{}shouts; cries out\textquotedblright{}.

\begin{quote}
\sk{सः} \sk{नरः} \sk{उच्चैः} \sk{क्रोशति।}
\end{quote}

\begin{quote}
\emph{saḥ naraḥ uccaiḥ krośati.}
\end{quote}

\begin{quote}
The man shouts loudly.
\end{quote}

Past (he …): \sk{सः} \sk{उच्चैः} \sk{क्रुष्टवान्}, \emph{saḥ uccaiḥ kruṣṭavān}.

\begin{cousinweb}
\textbf{\sk{क्रोशति}} (\emph{krośati}) comes from the root √kruś, \textquotedblleft{}to cry out, shout\textquotedblright{}; MW compares Greek \emph{krazō} (MW).
\end{cousinweb}

\begin{grammarlens}[title={What you've built}]
One more word for telling what happened.
\end{grammarlens}

\subsection*{Your turn}
\begin{itemize}
\raggedright
  \item \emph{Say it:} \emph{krośati}
  \item \emph{Say it:} \emph{krośati}, once more
  \item \emph{Say it:} \emph{saḥ naraḥ uccaiḥ krośati.} — The man shouts loudly.
  \item \emph{Recall:} say the Sanskrit for a heap, then the Sanskrit for in distress, then say \emph{krośati} again
\end{itemize}

\begin{tcolorbox}[breakable,colback=teal!4,colframe=teal!35!black,title={Before you move on}]
What does \sk{क्रोशति} mean? (\textquotedblleft{}Shouts; cries out\textquotedblright{}.) Say it once more.
\end{tcolorbox}
\section[\texorpdfstring{\sk{विलपति} (vilapati)}{विलपति (vilapati)}]{\sk{विलपति} (vilapati) — laments; wails}

\label{lesson:SA-C304-vilapati}



\begin{quote}
Before the new one: say the Sanskrit for in distress, then the Sanskrit for shouts.
\end{quote}

\subsection*{You'll want to know: \sk{विलपति}}
\textbf{\sk{विलपति}} — \emph{vilapati} — \textquotedblleft{}laments; wails\textquotedblright{}.

\begin{quote}
\sk{बालस्य} \sk{माता} \sk{विलपति।}
\end{quote}

\begin{quote}
\emph{bālasya mātā vilapati.}
\end{quote}

\begin{quote}
The boy's mother wails.
\end{quote}

Past (she …): \sk{सा} \sk{बहु} \sk{विलपितवती}, \emph{sā bahu vilapitavatī}.

\begin{cousinweb}
\textbf{\sk{विलपति}} (\emph{vilapati}) is built from \emph{vi} and the root √lap, \textquotedblleft{}to talk, prate\textquotedblright{}; MW gives \emph{vi-lap} as to utter moaning sounds, wail, lament (MW).
\end{cousinweb}

\begin{grammarlens}[title={What you've built}]
One more word for telling what happened.
\end{grammarlens}

\subsection*{Your turn}
\begin{itemize}
\raggedright
  \item \emph{Say it:} \emph{vilapati}
  \item \emph{Say it:} \emph{vilapati}, once more
  \item \emph{Say it:} \emph{bālasya mātā vilapati.} — The boy's mother wails.
  \item \emph{Recall:} say the Sanskrit for in distress, then the Sanskrit for shouts, then say \emph{vilapati} again
\end{itemize}

\begin{tcolorbox}[breakable,colback=teal!4,colframe=teal!35!black,title={Before you move on}]
What does \sk{विलपति} mean? (\textquotedblleft{}Laments; wails\textquotedblright{}.) Say it once more.
\end{tcolorbox}
\section[\texorpdfstring{\sk{निःसारयति} (niḥsārayati)}{निःसारयति (niḥsārayati)}]{\sk{निःसारयति} (niḥsārayati) — brings out; gets out}

\label{lesson:SA-C304-nihsarayati}



\begin{quote}
Before the new one: say the Sanskrit for shouts, then the Sanskrit for laments.
\end{quote}

\subsection*{You'll want to know: \sk{निःसारयति}}
\textbf{\sk{निःसारयति}} — \emph{niḥsārayati} — \textquotedblleft{}brings out; gets out\textquotedblright{}.

\begin{quote}
\sk{जनाः} \sk{तं} \sk{नरं} \sk{राशेः} \sk{निःसारयन्ति।}
\end{quote}

\begin{quote}
\emph{janāḥ taṃ naraṃ rāśeḥ niḥsārayanti.}
\end{quote}

\begin{quote}
People bring the man out of the heap.
\end{quote}

Past (he …): \sk{सः} \sk{नरं} \sk{निःसारितवान्}, \emph{saḥ naraṃ niḥsāritavān}.

\begin{cousinweb}
\textbf{\sk{निःसारयति}} (\emph{niḥsārayati}) is built from \emph{niḥ}, \textquotedblleft{}out\textquotedblright{}, and the causative of the root √sṛ, \textquotedblleft{}to go, run\textquotedblright{}; MW gives \emph{niḥ-sārayati} as to cause to go out, turn out (MW).
\end{cousinweb}

\begin{grammarlens}[title={What you've built}]
One more word for telling what happened.
\end{grammarlens}

\subsection*{Your turn}
\begin{itemize}
\raggedright
  \item \emph{Say it:} \emph{niḥsārayati}
  \item \emph{Say it:} \emph{niḥsārayati}, once more
  \item \emph{Say it:} \emph{janāḥ taṃ naraṃ rāśeḥ niḥsārayanti.} — People bring the man out of the heap.
  \item \emph{Recall:} say the Sanskrit for shouts, then the Sanskrit for laments, then say \emph{niḥsārayati} again
\end{itemize}

\begin{tcolorbox}[breakable,colback=teal!4,colframe=teal!35!black,title={Before you move on}]
What does \sk{निःसारयति} mean? (\textquotedblleft{}Brings out; gets out\textquotedblright{}.) Say it once more.
\end{tcolorbox}
